% =====================================================================
%  CSE 402 - Numerical Analysis, Simulation and Modeling Sessional
%  Final Project Report
%
%  Template: ACM Conference Proceedings (acmart, sigconf).
%  Build:    pdflatex main.tex  (twice; the bibliography is inline)
%
%  Every number in this report is read from a table under results/ in the
%  repository, and every figure is drawn by scripts/make_report_figures.py.
%  The README maps each table and figure to the script that produces it.
% =====================================================================
\documentclass[sigconf,nonacm]{acmart}

\settopmatter{printacmref=false}
\renewcommand\footnotetextcopyrightpermission[1]{}
\pagestyle{plain}

\usepackage{booktabs}
\usepackage{amsmath}
\usepackage{graphicx}

\setlength{\tabcolsep}{4pt}

% Figures live in the repository's figures/ folder, one level up.
\newcommand{\figdir}{../figures/}
\newcommand{\fig}[2][\linewidth]{%
  \IfFileExists{\figdir#2}{\includegraphics[width=#1]{\figdir#2}}{%
    \fbox{\parbox{0.85\linewidth}{\footnotesize\centering
      \texttt{\detokenize{#2}} is missing.\\
      Run \texttt{python scripts/make\_report\_figures.py}.}}}}

\begin{document}

\title{Rooting for Kepler: Testing, Explaining and Repairing
With-Memory Root Finders on Kepler's Equation}

\subtitle{CSE 402 Final Project Report \textbar\ Section A, Group 1}

\author{Anisa Binte Asad}
\affiliation{\institution{Student ID 2105036, BUET}\country{Bangladesh}}
\author{Dipit Saha}
\affiliation{\institution{Student ID 2105050, BUET}\country{Bangladesh}}
\author{Atika Tabassum Suchi}
\affiliation{\institution{Student ID 2105053, BUET}\country{Bangladesh}}
\author{S. M. A. M. Mahdi}
\affiliation{\institution{Student ID 2105056, BUET}\country{Bangladesh}}
\author{Fariha Ifrat}
\affiliation{\institution{Student ID 2105059, BUET}\country{Bangladesh}}

\renewcommand{\shortauthors}{Section A, Group 1}

% =====================================================================
\begin{abstract}
Kepler's equation, $E - e\sin E = M$, has to be solved every time an
orbit is evaluated, and a radial-velocity fit evaluates orbits hundreds
of thousands of times. Two recent papers propose with-memory iterative
methods, NWM9 (2024) and NWM11 (2025), with convergence orders $8.90$ and
$10.74$ and no extra function evaluations. Neither paper tests them on
Kepler's equation. We reimplemented both, confirmed their orders at
2000-digit precision, and compared them with Newton's, Danby's and
Markley's methods over 125{,}800 solves on a 7{,}400-point grid of
eccentricity $e$ and mean anomaly $M$, with four different starting
guesses. The new methods need the fewest iterations (1.38 for NWM11
against 1.94 for Danby), but near $e \to 1$, $M \to 0$ they fail to
converge 33 to 37 times as often as Danby. We traced these failures to
the self-accelerating parameter: where $f'(E) = 1 - e\cos E$ is close to
zero, the parameter's correction term becomes as large as the derivative
it is meant to adjust. We propose a one-line safeguard that keeps the
parameter only when its term is less than half of $|f'|$. Safeguarded
NWM11 cuts its failures from 74 to 3 (Danby has 2), keeps its measured
order of $10.74$ up to $e = 0.999$, and needs the same number of
iterations. We also find that the starting guess matters as much as the
method: Napier's guess lowers NWM9's mean iteration count from 2.37 to
1.44. Finally, solver error shifts the fitted orbit at least 1{,}000 times
less than measurement noise does.
\end{abstract}

\keywords{Kepler's equation, with-memory iterative methods, convergence
order, root finding, safeguarding, radial velocity, benchmarking}

\maketitle

% =====================================================================
\section{Introduction and Base Paper Review}

\subsection{The problem}

A planet's position along its orbit is described by the mean anomaly $M$,
which grows linearly with time, and the eccentric anomaly $E$, which fixes
where the planet actually is. The two are linked by Kepler's equation
\begin{equation}
f(E) = E - e\sin E - M = 0, \qquad 0 \le e < 1,
\label{eq:kepler}
\end{equation}
where $e$ is the orbital eccentricity. The equation has no closed-form
solution, so it is solved numerically every time an orbit is evaluated.
RadVel \cite{radvel}, a standard toolkit for fitting exoplanet orbits to
radial-velocity data, solves it once per observation for every
likelihood evaluation. In our experiments a single RadVel fit needed about
$1.3\times10^{5}$ solves (Section~\ref{sec:downstream}), and a Markov chain
Monte Carlo run multiplies that by thousands. The speed and reliability of
the solver therefore matter in practice.

Two features of Equation~\eqref{eq:kepler} shape everything that follows.
First, its derivative $f'(E) = 1 - e\cos E$ approaches zero as $e \to 1$
and $M \to 0$, because the root $E$ is then close to zero and
$\cos E \approx 1$. Every Newton-type step divides by $f'$, so this
corner of the $(e, M)$ plane is where solvers struggle. Second, $f$ and
$f'$ are built from $\sin E$ and $\cos E$ at the same point, so a method
that needs both pays for one pair of trigonometric calls, not two separate
evaluations. A method's cost on this equation is therefore not the same as
its cost on a generic nonlinear equation.

RadVel's own solver is Danby's fourth-order method \cite{danby}, written in
C following Murray and Dermott \cite{murray}, with a residual tolerance of
$10^{-12}$ \cite{radvelsrc}. That is the method a new solver has to beat.

\subsection{The base papers}

Our base papers are two articles by Mittal, Panday and co-authors
\cite{nwm9,nwm11}. Both start from an optimal eighth-order, three-step
method and raise its order with \emph{self-accelerating parameters}:
quantities that are recomputed every iteration from points the method has
already evaluated. Because the method reuses information from earlier
iterations, it is called a \emph{with-memory} method.

\textbf{NWM9} \cite{nwm9} adds one parameter, $\alpha_n$, to the third
step of the eighth-order method of Matthies et al.\ \cite{matthies}. The
paper proves an R-order of $8.8989$ and an efficiency index that rises from
$1.6818$ to $1.7272$.

\textbf{NWM11} \cite{nwm11} adds two parameters, $\alpha_k$ and $\beta_k$,
to the first and third steps of the eighth-order method of Solaiman and
Hashim \cite{solaiman}. The paper proves an R-order of $10.7446$ and an
efficiency index of $1.8105$.

In both methods the parameters are derivatives of Hermite interpolating
polynomials fitted through the current and previous iterates. Those points
have already been evaluated, so the higher order costs no new function
evaluations. That is the central claim of both papers.

\subsection{What the base papers leave open}

Both papers test their methods on eight smooth scalar functions, each from
a single starting point, and measure cost with the efficiency index
$\rho^{1/\gamma}$, where $\rho$ is the order and $\gamma$ counts function
evaluations. This leaves four questions open for Kepler's equation.

\begin{enumerate}
\item \textbf{Cost.} The efficiency index counts every evaluation as equal.
      On Kepler's equation a $(\sin, \cos)$ pair, a lone $\sin$ and an
      interpolated derivative all cost different amounts.
\item \textbf{Behaviour across the domain.} Eight test functions say
      nothing about how a method behaves across the whole $(e, M)$ plane,
      and in particular near $e \to 1$, $M \to 0$ where $f' \to 0$.
\item \textbf{Robustness.} Neither paper reports how often the methods
      fail to converge, or why.
\item \textbf{Relevance.} Neither paper asks whether a higher order changes
      anything in an application whose data are noisy.
\end{enumerate}

\subsection{Our contributions}

We keep the published formulas exactly as given and build the evaluation
the papers do not have. Then, having found where the methods break, we
modify them. Our contributions are:

\begin{enumerate}
\item \textbf{A verified reimplementation.} Both methods reproduce their
      published orders at 2000-digit precision. We also show that matching
      the order is not enough to prove a correct transcription, and that
      each paper's error constant is the check that works
      (Section~\ref{sec:verify}).
\item \textbf{A Kepler-specific benchmark} of five solvers and four
      starting guesses over 125{,}800 solves, with cost counted in measured
      units of trigonometric work (Sections~\ref{sec:grid} and
      \ref{sec:guess}).
\item \textbf{An explanation of the failures.} The with-memory methods fail
      in the corner because of the self-accelerating parameter, not the
      underlying eighth-order method. We show this directly by switching
      the memory off (Section~\ref{sec:failure}).
\item \textbf{An improvement to the base method.} A safeguard on the
      parameter that removes almost all of NWM11's failures without
      changing its order or its iteration count (Section~\ref{sec:fix}).
\item \textbf{A downstream check} showing how far solver error travels
      into a fitted orbit, compared with the noise in real data
      (Section~\ref{sec:downstream}).
\end{enumerate}

All code, configuration files, result tables and figure scripts are public
at \url{https://github.com/kittocatto1/CSE402-Rooting-for-Kepler-}.

% =====================================================================
\section{Methodology}

\subsection{The with-memory iteration}
\label{sec:nwm}

We use NWM11 to show the structure; NWM9 is built the same way with one
parameter instead of two. Starting from $s_k$, one NWM11 iteration is
\begin{align}
v_k     &= s_k - \frac{f(s_k)}{f'(s_k) + \alpha_k f(s_k)}, \label{eq:step1}\\
t_k     &= v_k - \frac{f(v_k)}{q}
           - \frac{2 f(v_k)^2 q R}{4q^4 - 4f(v_k)q^2R + f(v_k)^2R^2}, \notag\\
s_{k+1} &= t_k - \frac{f(t_k)}{w' + \beta_k f(t_k)}, \label{eq:step3}
\end{align}
where $q$, $R$ and $w'$ are divided-difference estimates of $f'(v_k)$,
$f''(v_k)$ and $f'(t_k)$ built from values already in hand
\cite{nwm11}. The parameters are
$\alpha_k = -H_5''(s_k)/(2f'(s_k))$ and
$\beta_k = H_7''''(t_k)/(4H_6'''(v_k)) + \alpha_k$, where $H_5$, $H_6$ and
$H_7$ are Hermite polynomials through the current and previous iterates.
With $\alpha_k = \beta_k = 0$ the scheme is the eighth-order method of
Solaiman and Hashim. Each iteration evaluates $f$ at three points but
needs $f'$ only at $s_k$, so it uses one $(\sin, \cos)$ pair and two lone
$\sin$ calls.

Two implementation details decide whether the parameters come out right.
The node lists repeat a point, for example
$[s_k, s_k, t_{k-1}, v_{k-1}, s_{k-1}, s_{k-1}]$, which makes an ordinary
divided difference $0/0$. Our interpolation code therefore uses the
confluent (Hermite) form and takes $f'$ at a repeated node. The papers
also need the derivative $H^{(k)}$, not the Newton coefficient $H^{(k)}/k!$.
Using the coefficient is an error of a factor $k!$ that still converges,
so it would not show up as a crash.

\subsection{Measuring convergence order}
\label{sec:coc}

We check each claimed order instead of assuming it, using the
computational order of convergence
\begin{equation}
p \approx \frac{\ln|\varepsilon_{n+1}/\varepsilon_{n}|}
               {\ln|\varepsilon_{n}/\varepsilon_{n-1}|},
\qquad \varepsilon_n = |x_n - x^{*}|,
\end{equation}
which needs three consecutive errors that are still above the precision
floor. A method of order 9 starting at an error of $10^{-1}$ reaches
$10^{-9}$, then $10^{-81}$, then $10^{-729}$. In double precision only two
of those errors are visible, which is not enough for one estimate. We
therefore measure order at 2000 decimal digits with \texttt{mpmath}
\cite{mpmath}; the same solver code runs on both floats and
\texttt{mpmath} numbers. Everywhere else the solvers run in ordinary
double precision.

\subsection{Counting cost on Kepler's equation}
\label{sec:costmodel}

Counting iterations is unfair, because one NWM11 iteration does far more
work than one Newton iteration. Counting evaluations equally, as the
efficiency index does, is also unfair here, because it charges a
$(\sin,\cos)$ pair the same as a lone $\sin$. Every solver in our code
evaluates $f$ through a shared problem object that records what each call
needed. We then price each kind of call in units of one \texttt{sin} call,
timed on the benchmark machine (Table~\ref{tab:weights}).

\begin{table}[h]
\centering\small
\caption{Cost of each primitive, in units of one \texttt{sin} call. The
measured weights are the median of three timing runs.}
\label{tab:weights}
\begin{tabular}{@{}lrr@{}}
\toprule
\textbf{Primitive} & \textbf{Measured (Python)} & \textbf{Call count} \\
\midrule
$\sin$ alone                  & 1.00 & 1 \\
$\cos$ alone                  & 1.04 & 1 \\
$(\sin,\cos)$ pair            & 2.34 & 1 \\
Interpolated derivative       & 1.76 & 0 \\
\bottomrule
\end{tabular}
\end{table}

We report two costings. The \emph{measured} weights are what our Python
code actually pays. The \emph{call-count} weights describe a compiled
implementation with a fused \texttt{sincos} instruction, where a pair costs
one call and arithmetic is nearly free; this is the accounting the base
papers' efficiency index assumes. The choice between them changes how the
classical and with-memory methods compare, so we show both.

\subsection{A safeguard for the self-accelerating parameter}
\label{sec:guard}

Each parameter enters a denominator as a correction added to a
derivative: $f'(s_k) + \alpha_k f(s_k)$ in Equation~\eqref{eq:step1}, and
$w' + \beta_k f(t_k)$ in Equation~\eqref{eq:step3}. The papers derive the
parameters assuming the iterate is already close to the root. There $f$ is
tiny and the correction is a small adjustment. Far from the root, or where
$f'$ is itself close to zero, the correction can be as large as the
derivative it is meant to adjust, and the step is thrown far off.

Our safeguard tests this directly. In each iteration, a parameter $\lambda$
(standing for $\alpha_k$, $\beta_k$ or NWM9's $\alpha_n$) is used only if
\begin{equation}
|\lambda\, f| \le c\, |f'(s_k)|, \qquad c = \tfrac{1}{2},
\label{eq:guard}
\end{equation}
where $f$ is the function value the parameter multiplies and $f'(s_k)$ is
the one derivative the iteration actually evaluates. If the test fails,
that parameter is set to zero for that iteration, which gives the step of
the memoryless eighth-order method. The test costs two multiplications and
no new evaluations.

The safeguard cannot lower the order. Near a simple root, $f \to 0$ while
$f'(s_k) \to f'(E^*) \ne 0$, and the parameters stay bounded, so
$|\lambda f|$ eventually falls below $c|f'|$ and the test always passes.
From then on the guarded method takes exactly the published steps, so it
has the published asymptotic order. We confirm this numerically in
Section~\ref{sec:fix}. We chose $c = 1/2$ after trying $0.5$, $0.1$ and
$0.01$ on the grid: all three removed most failures, and $0.5$ removed the
most while changing the fewest steps.

\subsection{Starting guesses}
\label{sec:guesses}

Every iterative method needs a first estimate $E_0$. We treat the guess as
a separate factor and test four:
$E_0 = M$ (\emph{simple});
$E_0 = M + e\sin M$ (\emph{canonical}, the first term of the series
solution);
$E_0 = M + 0.85\,e\,\mathrm{sign}(\sin M)$ (\emph{radvel}, the guess in
RadVel's code \cite{radvelsrc}); and
$E_0 = e\sin M + \max\{M,\ e(\sin M + 0.591)\}$ (\emph{napier}, found by
symbolic regression \cite{napier}).
Markley's method \cite{markley} is a closed-form starter followed by one
fixed correction, so it takes no guess and runs once per point.

% =====================================================================
\section{Implementation and Experiments}

\subsection{Code structure}

The code is a Python package, \texttt{keplerbench}. Every solver and every
starting guess implements one shared interface, so the experiment runner
treats them all the same way and every iterative solver uses the same
stopping rule: stop when $|f(E)| \le 10^{-14}$, or give up after 50
iterations. Solvers never call \texttt{math.sin} directly. They evaluate
Equation~\eqref{eq:kepler} through a \texttt{KeplerProblem} object that
counts every $(\sin,\cos)$ pair, lone $\sin$ or $\cos$, and interpolated
derivative. Because the counting happens outside the solver, a solver
cannot under-report its own cost. Table~\ref{tab:solvers} lists the
solvers.

\begin{table}[h]
\centering\small
\caption{The solvers under test.}
\label{tab:solvers}
\begin{tabular}{@{}llr@{}}
\toprule
\textbf{Method} & \textbf{Type} & \textbf{Order} \\
\midrule
Newton--Raphson                & classical            & 2 \\
Danby \cite{danby}             & classical            & 4 \\
Markley \cite{markley}         & closed form          & --- \\
NWM9 \cite{nwm9}               & with memory          & 8.8989 \\
NWM11 \cite{nwm11}             & with memory          & 10.7446 \\
NWM9, NWM11 guarded (ours)     & with memory, guarded & as published \\
\bottomrule
\end{tabular}
\end{table}

\subsection{Reference roots}

Accuracy is measured against roots computed separately in \texttt{mpmath}
at 50 decimal digits, by bisection followed by high-precision Newton steps,
with residuals checked to be below $10^{-40}$. This code shares nothing
with the solvers under test, so no solver is graded against itself.

\subsection{The test grid}
\label{sec:gridsetup}

We test on 7{,}400 $(e, M)$ points drawn in three different ways
(Table~\ref{tab:grid-sets}). The even grid covers the whole domain fairly.
The corner set is log-spaced, because equal steps would put almost no
points near $e \to 1$, $M \to 0$. The RadVel set follows the eccentricity
distribution of real exoplanets \cite{kipping}. The grid stops at
$M = \pi$ because Equation~\eqref{eq:kepler} is symmetric about it.

\begin{table}[h]
\centering\small
\caption{The three sets of test points.}
\label{tab:grid-sets}
\begin{tabular}{@{}lllr@{}}
\toprule
\textbf{Set} & \textbf{$e$} & \textbf{$M$} & \textbf{Points} \\
\midrule
Even grid  & 50 values, 0 to 0.99         & 100 values, 0 to $\pi$        & 5{,}000 \\
Corner set & $1-e$ from $10^{-1}$ to $10^{-4}$ & $10^{-1}$ to $10^{-6}$ & 400 \\
RadVel set & Beta(0.867, 3.03)            & uniform on $[0, 2\pi)$         & 2{,}000 \\
\bottomrule
\end{tabular}
\end{table}

When we report failure rates we use one fixed definition of the
\emph{hard corner}: $e > 0.9$ and $M < 0.1$. It holds 381 of the 7{,}400
points. Crossing four iterative solvers with four guesses, plus Markley
once, gives 17 combinations and $17 \times 7{,}400 = 125{,}800$ solves. Each
iterative solver makes $4 \times 381 = 1{,}524$ solves in the hard corner.

\subsection{Experiments}

Table~\ref{tab:experiments} lists the experiments. Each one is driven by a
configuration file and writes its results, with a record of the git commit
and configuration that produced them, under \texttt{results/}.

\begin{table}[h]
\centering\small
\caption{Experiments, their scripts and their outputs.}
\label{tab:experiments}
\begin{tabular}{@{}p{1.55cm}p{3.7cm}p{2.3cm}@{}}
\toprule
\textbf{Experiment} & \textbf{Question} & \textbf{Output} \\
\midrule
verification   & Does each solver reproduce its paper's order, and find the
                 right root? & \texttt{results/ verification/} \\
grid benchmark & Iterations, cost, time and failures over 125{,}800
                 solves & \texttt{results/ grid\_benchmark/} \\
safeguard      & Do the failures come from the memory, and does the
                 safeguard remove them? & \texttt{results/ safeguard/} \\
propagation    & How far does solver error move a fitted orbit?
               & \texttt{results/ error\_propagation/} \\
Monte Carlo    & How far does measurement noise move it?
               & \texttt{results/ monte\_carlo/} \\
real data      & Do the solvers agree on published data?
               & \texttt{results/ real\_data/} \\
\bottomrule
\end{tabular}
\end{table}

The downstream experiments use the real observation times, uncertainties
and three-instrument structure of HD~164922 \cite{hd164922} (401
measurements), with the velocities replaced by a known orbit
($P = 20.885$\,d, $e = 0.35$, $\omega = 1.0$\,rad, $K = 5$\,m\,s$^{-1}$)
plus Gaussian noise at the real uncertainties. A known orbit lets us
measure how far each fit lands from the truth. Each solver is patched into
RadVel in place of its compiled solver, and a call counter checks after
every fit that the patch was actually used.

\subsection{Reproducibility}

All experiments and figures were rerun on the current code for this
report, and each result table records the commit that produced it. The repository's README maps each table and figure to the script
and result file behind it. The test suite has 284 tests, including checks
of every solver against its paper and of the safeguard's three claims.

% =====================================================================
\section{Results and Discussion}

\subsection{Verification}
\label{sec:verify}

Table~\ref{tab:order} shows the measured order of each iterative solver on
its own paper's eight test functions, at 2000 digits. All four pass. The
mean measured order is at or slightly above each claim, which is expected
because the papers state R-order as a lower bound. Single functions scatter
by a fraction of a percent on either side (Figure~\ref{fig:order}).

\begin{table}[h]
\centering\small
\caption{Measured convergence order on each paper's eight test functions.}
\label{tab:order}
\begin{tabular}{@{}lrrl@{}}
\toprule
\textbf{Solver} & \textbf{Claimed} & \textbf{Measured} & \textbf{Range} \\
\midrule
Newton & 2.0000  & 2.0000  & 2.0000--2.0000 \\
Danby  & 4.0000  & 4.0000  & 4.0000--4.0000 \\
NWM9   & 8.8989  & 8.9262  & 8.8851--9.0000 \\
NWM11  & 10.7446 & 10.7605 & 10.7147--10.8336 \\
\bottomrule
\end{tabular}
\end{table}

\begin{figure}[h]
\centering
\fig{convergence/measured_vs_claimed_order.pdf}
\caption{Measured order against each paper's claim. Error bars show the
spread over the eight test functions.}
\label{fig:order}
\end{figure}

\paragraph{Matching the order does not prove the code is right.} Both
papers print nested fractions that can be read in more than one way, and
several wrong readings still converge at order 8. Measuring the order
would not have caught them. Each paper's eighth-order \emph{error
constant} did. We drove each scheme one step from a point $10^{-24}$ away
from the root and compared the measured constant with the published
formula; our implementations agree to better than $10^{-20}$ in relative
terms. One wrong grouping in NWM11's $w'$ missed the constant by about 60
orders of magnitude while still converging at order 8. The same check
caught test functions that had been mis-copied from the PDFs, where
superscripts had detached during text extraction (six in the NWM9 paper).
As a further check, setting every self-accelerating parameter to zero
gives exactly order $8.0$ on all eight functions for both methods, which
confirms that the extra order comes from the memory.

\paragraph{The right root, not just a small residual.} A solver can
converge to a wrong branch and still report a tiny residual. We ran each
solver on 4{,}975 points against the 50-digit reference
(Table~\ref{tab:correct}). No solver returned a wrong root anywhere. The
errors of order $10^{-10}$ come from points where $f'$ is small: the
stopping rule tests the residual, and the error in $E$ is roughly
$f/f'$. Where a solver does not converge it uses up its 50 iterations and
says so. That outcome is visible and can be retried, unlike a wrong root,
so we treat it as a robustness problem in Section~\ref{sec:failure}.

\begin{table}[h]
\centering\small
\caption{Correctness on 4{,}975 points against the 50-digit reference
(stopping tolerance $10^{-13}$).}
\label{tab:correct}
\begin{tabular}{@{}lrrr@{}}
\toprule
\textbf{Solver} & \textbf{Converged} & \textbf{Wrong roots} & \textbf{Largest $|E-E_{\mathrm{ref}}|$} \\
\midrule
Newton & 99.54\% & 0 & $3.4\times10^{-10}$ \\
Danby  & 99.72\% & 0 & $1.6\times10^{-10}$ \\
NWM9   & 98.25\% & 0 & $2.3\times10^{-10}$ \\
NWM11  & 99.56\% & 0 & $2.2\times10^{-10}$ \\
\bottomrule
\end{tabular}
\end{table}

\subsection{Iterations, cost and time on the grid}
\label{sec:grid}

Table~\ref{tab:grid} shows each solver with the starting guess that suits
it best. The three columns tell three different stories.

\begin{table}[h]
\centering\small
\caption{Each solver with its best guess, over all 7{,}400 points. Cost is
in measured \texttt{sin} units; time is the median per solve.}
\label{tab:grid}
\begin{tabular}{@{}llrrrr@{}}
\toprule
\textbf{Solver} & \textbf{Guess} & \textbf{Mean} & \textbf{Median} & \textbf{Time} & \textbf{Conv.} \\
 & & \textbf{iter.} & \textbf{cost} & \textbf{($\mu$s)} & \\
\midrule
Markley & ---       & 0.00 & 3.34  & 1.7 & 100.00\% \\
NWM11   & canonical & 1.38 & 6.34  & 2.8 & 99.93\% \\
NWM9    & napier    & 1.44 & 6.34  & 3.7 & 100.00\% \\
Danby   & napier    & 1.94 & 7.68  & 2.0 & 100.00\% \\
Newton  & napier    & 3.21 & 11.02 & 2.1 & 100.00\% \\
\bottomrule
\end{tabular}
\end{table}

\paragraph{Iterations.} The with-memory methods need the fewest
iterations. The median solve takes one iteration for NWM9 and NWM11, two
for Danby and three for Newton. The higher order is real and visible.

\paragraph{Cost.} In measured units, the median with-memory solve costs
$6.34$ against Danby's $7.68$, about 17\% less. The arithmetic can be
checked by hand. Danby's typical solve is two iterations: two
$(\sin,\cos)$ pairs and three lone $\sin$ calls (one of them the stopping
test), so $2(2.34) + 3 = 7.68$. The typical with-memory solve is a single
iteration: one pair and four lone $\sin$ calls, so $2.34 + 4 = 6.34$. This
saving needs care. In a one-iteration solve there is no previous iterate,
so the parameters are still at their fixed starting values and no memory
is used. The 17\% saving therefore belongs to the eighth-order parent
methods, which is also why NWM9 and NWM11 tie. The memory is used only in
longer solves, and there a with-memory iteration costs $10.62$ against
$3.34$ for a classical one, because each one adds three interpolated
derivatives at $1.76$ each (Figure~\ref{fig:costbd}). Under call counting
the same two iterations cost $4$ and $2$.

\begin{figure*}[t]
\centering
\fig[0.85\textwidth]{cost/cost_breakdown.pdf}
\caption{Cost of one iteration. Left: weights measured in our Python
code. Right: call counting, as in a compiled code with a fused
\texttt{sincos}. The difference between the two panels is almost entirely
the interpolated derivatives, which cost $1.76$ each in Python and close
to nothing in compiled code.}
\label{fig:costbd}
\end{figure*}

\paragraph{Time.} On the clock the ranking reverses. Markley, Danby and
Newton take 1.7 to 2.1\,$\mu$s per solve. NWM11 at its best guess takes
2.8\,$\mu$s and NWM9 3.7\,$\mu$s, 1.4 and 1.8 times as long as Danby, and
NWM11 under the simple guess takes 9.5\,$\mu$s. The model says the
with-memory methods are cheaper; the clock says they are slower.
Section~\ref{sec:sincos} explains why.

\subsection{The starting guess}
\label{sec:guess}

The base papers say nothing about the starting guess, yet for NWM9 it
matters more than the choice of method. Figure~\ref{fig:byguess} shows the
mean iteration count of every solver under every guess. Napier's guess is
the best for Newton, Danby and NWM9, and the canonical guess is best for
NWM11, with Napier's close behind.

\begin{figure}[h]
\centering
\fig{grid/iterations_by_guess.pdf}
\caption{Mean iterations to reach $|f| \le 10^{-14}$ for each solver and
guess, over all 7{,}400 points. The hatched bar is the solver's mean over
all four guesses.}
\label{fig:byguess}
\end{figure}

Changing NWM9's guess from $E_0 = M$ to Napier's formula lowers its mean
from $2.37$ to $1.44$ iterations, a saving of $0.93$. Moving from Danby to
NWM9 with Napier's guess saves only $0.51$ ($1.94 \to 1.44$). For Newton the
order is reversed: the better guess saves $0.64$ ($3.84 \to 3.21$), while
switching to Danby saves $1.26$. The guess also changes reliability. NWM9
converges on $99.43\%$ of points from the simple guess and on all of them
from Napier's.

Splitting the averages by set shows that the sets do not agree
(Figure~\ref{fig:byset}). On the even grid Napier's guess is the best for
every solver. On the RadVel set, Napier's and the canonical guess are
within $0.01$ iterations of each other, and both beat $E_0 = M$. In the
corner set no single guess wins: Napier's guess rescues NWM9, cutting its
mean from $7.05$ to $2.77$, while NWM11 does best from $E_0 = M$.

\begin{figure*}[t]
\centering
\fig[0.86\textwidth]{grid/iterations_by_set.pdf}
\caption{Mean iterations in each set of test points. Orange marks the
fewest iterations for that solver in that set. The mean under each solver
is its average over all four guesses. Only converged solves are counted.}
\label{fig:byset}
\end{figure*}

Figure~\ref{fig:heat} shows where the iterations are spent. Over most of
the domain every solver needs one to three iterations. The extra work is
concentrated at high $e$ and low $M$.

\begin{figure*}[t]
\centering
\fig[0.78\textwidth]{grid/iterations_best_guess.pdf}
\caption{Iterations to tolerance on the even grid, each solver with its
best guess. All four solvers are cheap almost everywhere; the cost is
concentrated near $e \to 1$, $M \to 0$.}
\label{fig:heat}
\end{figure*}

% ---------------------------------------------------------------------
\subsection{Where and why the with-memory methods fail}
\label{sec:failure}

\subsubsection{Failures concentrate in the hard corner}

Table~\ref{tab:robust} splits non-convergence by region. Of 176 failures
in 125{,}800 solves, 155 are in the hard corner, which holds only 5\% of
the solves. Every failure is a solve that used up its 50 iterations. No
solver crashed or returned a wrong root.

\begin{table}[h]
\centering\small
\caption{Non-convergence by region, over all four guesses. The hard corner
is $e > 0.9$, $M < 0.1$ (1{,}524 solves per iterative solver).}
\label{tab:robust}
\begin{tabular}{@{}lrrr@{}}
\toprule
\textbf{Solver} & \textbf{Rest of grid} & \textbf{Hard corner} & \textbf{Failures} \\
\midrule
Markley & 0.000\% & 0.000\% & 0 \\
Danby   & 0.000\% & 0.131\% & 2 \\
Newton  & 0.018\% & 0.919\% & 19 \\
NWM9    & 0.057\% & 4.265\% & 81 \\
NWM11   & 0.000\% & 4.856\% & 74 \\
\bottomrule
\end{tabular}
\end{table}

Outside the corner NWM11 never fails. Inside it, NWM11 fails on 4.9\% of
solves, 37 times Danby's rate, and NWM9 on 4.3\%, 33 times Danby's rate.
Figure~\ref{fig:failmaps} shows where each solver fails.

\begin{figure*}[t]
\centering
\fig[0.9\textwidth]{grid/failure_maps_by_solver.pdf}
\caption{Converged and failed solves for each solver, all guesses
together. The vertical axis is $1-e$ on a log scale, so the top of each
panel is $e \to 1$; dashed lines mark the hard corner. Almost every
failure lies inside it.}
\label{fig:failmaps}
\end{figure*}

\subsubsection{The failures come from the memory}

As $e \to 1$ and $M \to 0$, the root $E$ goes to zero, $\cos E \to 1$, and
$f'(E) = 1 - e\cos E$ goes to zero. At $e = 0.9999$, $f'(0) = 10^{-4}$.
A Newton-type step divides by this derivative, so a small error in the
iterate becomes a large jump. All Newton-type methods feel this, which is
why Newton and even Danby have a few failures here. The with-memory
methods have a second problem on top of it. Their parameter is added to a
denominator built from $f'$, as in $f'(s_k) + \alpha_k f(s_k)$. The
parameter is itself computed from interpolated derivatives, and NWM9's
$\alpha_n$ depends on $f'$ in both of its terms, one of which divides by
it. When $f'$ is tiny, the correction $\alpha_k f(s_k)$ can be as large as
$f'$ itself. The denominator then no longer resembles the derivative it is
meant to approximate, and the step can be thrown far from the root.

We tested this explanation directly by running both methods with the
memory switched off: every parameter fixed at zero, which gives the
eighth-order parent methods. Table~\ref{tab:fix} shows the result.
NWM11's failures drop from 74 to 9 and NWM9's from 81 to 18, while the
mean iteration count barely changes. Most of the failures are therefore
caused by the memory, not by the base method.

\subsubsection{NWM9's order collapses at high eccentricity}
\label{sec:orderdrop}

The same effect shows up in the convergence order measured on Kepler's
equation itself (Table~\ref{tab:ordere}, Figure~\ref{fig:ordere}). NWM11
keeps its order of about $10.7$ all the way to $e = 0.999$. When it
converges in the corner, it converges at full order. NWM9 drops to $4.5$
at $e = 0.99$ and to $0.36$ at $e = 0.999$, which is slower than linear.
RadVel limits $e$ to $0.99$ \cite{radvelsrc}, so the worst of this
collapse lies outside its working range. The value $2.82$ at $e = 0.9$
is probably an artefact of the estimator rather than the method, because
the points on either side are both close to $8.9$. Newton's $1.79$ at
$e = 0.999$ is the same kind of artefact.

\begin{table}[h]
\centering\footnotesize
\setlength{\tabcolsep}{3.2pt}
\caption{Order measured on Kepler's equation at $M = 0.3$, 2000 digits.
A dash means the estimate could not be resolved.}
\label{tab:ordere}
\begin{tabular}{@{}lrrrrrrr@{}}
\toprule
$e$ & 0.3 & 0.7 & 0.9 & 0.95 & 0.99 & 0.995 & 0.999 \\
\midrule
Newton & 2.00 & 2.00 & 2.00 & 2.00 & 2.13 & --- & 1.79 \\
Danby  & 4.00 & 4.00 & 4.00 & 4.00 & 4.00 & 4.00 & 4.00 \\
NWM9   & 8.90 & 8.90 & 2.82 & 8.90 & 4.48 & 0.72 & 0.36 \\
NWM11  & 10.72 & 10.76 & 10.88 & 10.84 & 10.73 & 10.74 & 10.74 \\
\bottomrule
\end{tabular}
\end{table}

\begin{figure}[h]
\centering
\fig{convergence/order_vs_eccentricity.pdf}
\caption{Measured order against eccentricity. Dotted lines mark each
method's claimed order; hollow markers are estimates that could not be
resolved.}
\label{fig:ordere}
\end{figure}

\subsubsection{The cost model's premise fails in Python}
\label{sec:sincos}

The cost argument for these methods assumes that one call returns both
$\sin E$ and $\cos E$ for about the price of one. In our measurements a
pair costs $2.34$, against $2.04$ for a $\sin$ and a $\cos$ timed
separately. The discount does not exist. The reason is that CPython has no
fused \texttt{sincos}: \texttt{math.sin} and \texttt{math.cos} are two
separate C calls. The one call that returns both,
\texttt{cmath.exp(1j*x)}, costs about six \texttt{sin} calls. Interpolated
derivatives are also expensive in Python, at $1.76$ \texttt{sin} units
each, because a few interpreted arithmetic operations cost more than one C
call. This is what turns a cheaper model cost into a slower clock time
(Figure~\ref{fig:wall}). In a compiled kernel such as RadVel's, both
effects would shrink, so the papers' reasoning may well hold there.

\begin{figure}[h]
\centering
\fig{cost/wallclock_vs_cost.pdf}
\caption{Wall-clock time against modelled cost, one point per timed
solve. Over all points $r = 0.99$, but that figure is carried by a few
very long solves (up to 50 iterations at $e > 0.93$). Over typical
solves (88\% of points, cost at most 15) the correlation falls to
$r = 0.42$.}
\label{fig:wall}
\end{figure}

% ---------------------------------------------------------------------
\subsection{The safeguard repairs NWM11}
\label{sec:fix}

Switching the memory off everywhere removes most failures but gives up the
order the papers prove. The safeguard of Section~\ref{sec:guard} keeps the
memory where it is safe. Table~\ref{tab:fix} and Figure~\ref{fig:fix}
compare the three versions of each method on the full grid.

\begin{table}[h]
\centering\small
\caption{Published, memory-off and guarded versions on all 7{,}400 points
and four guesses (29{,}600 solves each). Mean iterations count converged
solves only.}
\label{tab:fix}
\begin{tabular}{@{}lrrr@{}}
\toprule
\textbf{Version} & \textbf{Failures} & \textbf{Corner rate} & \textbf{Mean iter.} \\
\midrule
Danby (reference)     & 2  & 0.13\% & 2.04 \\
\midrule
NWM9 published        & 81 & 4.27\% & 1.76 \\
NWM9 memory off       & 18 & 0.85\% & 1.74 \\
NWM9 guarded          & 17 & 0.85\% & 1.75 \\
\midrule
NWM11 published       & 74 & 4.86\% & 1.52 \\
NWM11 memory off      & 9  & 0.59\% & 1.51 \\
NWM11 guarded         & \textbf{3}  & \textbf{0.20\%} & 1.53 \\
\bottomrule
\end{tabular}
\end{table}

\begin{figure}[h]
\centering
\fig{safeguard/safeguard_failures.pdf}
\caption{Failed solves in the hard corner (out of 1{,}524) for the
published, memory-off and guarded versions of each method, with Danby for
reference.}
\label{fig:fix}
\end{figure}

\paragraph{NWM11.} The guard cuts NWM11's failures from 74 to 3. That is
fewer than switching the memory off (9), and close to Danby (2). With the
canonical guess, guarded NWM11 converges on all 7{,}400 points in $1.38$
iterations on average, against Danby's best of $1.94$. Outside the hard
corner the guard rarely acts: guarded NWM11 takes the same number of
iterations as the published version on 99.6\% of those solves, and its
largest error there is unchanged at $5.7\times10^{-14}$. The guard also
keeps the order. At 2000 digits, guarded NWM11 measures $10.76$ on
the paper's test functions, the same as the published version, and the
same value at every eccentricity in Table~\ref{tab:ordere}, up to $10.74$
at $e = 0.999$.

\paragraph{NWM9.} The guard cuts NWM9's failures from 81 to 17, the same
as switching the memory off (18). The 17 that remain also fail without
any memory, so they belong to NWM9's eighth-order parent method, not to
its parameter. The guard keeps NWM9's order on its test functions
($8.93$), but it does not fix the order collapse of
Section~\ref{sec:orderdrop}. For NWM9 the best protection is the guess:
with Napier's guess, published NWM9 already converges everywhere.

The guard is the improvement we propose to the base papers: a test that
costs two multiplications per parameter, needs no new evaluations, and
leaves the method unchanged wherever the papers' analysis applies. On
Kepler's equation it brings NWM11 to Danby's robustness while keeping its
lower iteration count.

% ---------------------------------------------------------------------
\subsection{Does solver error matter downstream?}
\label{sec:downstream}

We traced solver error from $E$ through the true anomaly $\nu$ and the
radial velocity $v_r$ into the fitted orbit. The question is how large
the solver's effect is compared with the uncertainty that is already in
the data. Table~\ref{tab:budget} compares the largest shift any solver
caused, at any tolerance, with the spread caused by noise across 50 refits
with fresh noise, and with RadVel's own MCMC posterior width.

\begin{table}[h]
\centering\small
\caption{Error budget. Solver shift: the largest shift over all five
solvers and all tolerances from $10^{-4}$ to $10^{-14}$ (in every case
Newton at $10^{-4}$). Noise: standard deviation over 50 refits with fresh
noise. MCMC: RadVel posterior width.}
\label{tab:budget}
\begin{tabular}{@{}lrrrr@{}}
\toprule
\textbf{Param.} & \textbf{Solver} & \textbf{Noise} & \textbf{MCMC} & \textbf{Noise/solver} \\
\midrule
$P$ (d)  & $1.0\times10^{-8}$ & $5.0\times10^{-4}$  & $6.4\times10^{-4}$  & $51{,}000$ \\
$e$      & $1.1\times10^{-5}$ & $1.65\times10^{-2}$ & $1.70\times10^{-2}$ & $1{,}500$ \\
$\omega$ & $2.3\times10^{-5}$ & $5.03\times10^{-2}$ & $5.59\times10^{-2}$ & $2{,}200$ \\
$K$      & $3.5\times10^{-5}$ & $9.93\times10^{-2}$ & $1.01\times10^{-1}$ & $2{,}800$ \\
\bottomrule
\end{tabular}
\end{table}

Even at a tolerance of $10^{-4}$, far looser than anyone would use, the
shift caused by the solver is at least 1{,}500 times smaller than the
shift caused by noise. Only Newton at $10^{-4}$ shows a real solver
effect, about $6.5\times10^{-4}$ standard deviations in $e$, because a
second-order method stopped at that tolerance still has an error of about
$10^{-4}$ in $E$. Every other combination stays below about $10^{-5}$
standard deviations, and that floor does not shrink as the tolerance
tightens. A floor that ignores the tolerance comes from the optimiser's
own stopping precision, not from the solver, so below $10^{-4}$ we can
only bound the solver's effect, not measure it
(Figure~\ref{fig:tol}). The shift is exactly zero at the $10^{-14}$
reference by construction, and Markley, which has no tolerance, gives
identical fits at every setting.

The spread across 50 noisy refits agrees with RadVel's MCMC posterior
width to within 3\% for $e$ and $K$, 10\% for $\omega$ and 21\% for $P$.
This is an independent check on the noise study itself.

\begin{figure*}[t]
\centering
\fig[0.8\textwidth]{propagation/tolerance_vs_parameter_shift.pdf}
\caption{Parameter shift against solver tolerance, in units of the noise
spread. Only Newton at $10^{-4}$ rises clearly above the floor set by the
optimiser.}
\label{fig:tol}
\end{figure*}

\paragraph{Cost of a whole fit.} Table~\ref{tab:fit} shows the time of a
full RadVel fit with each solver. All of our solvers are 18 to 32 times
slower than RadVel's compiled one. Most of that gap comes from calling a
scalar Python solver once per point, where RadVel works on whole arrays,
so these numbers compare our solvers with each other, not with production
code. Newton, Danby and Markley take 0.9 to 1.1\,s; the with-memory
methods are the slowest, by 1.3 to 1.7 times.

\begin{table}[h]
\centering\small
\caption{Time of one full RadVel fit (about $1.3\times10^{5}$ solves),
median over six tolerances. Each fit was timed once.}
\label{tab:fit}
\begin{tabular}{@{}lrrr@{}}
\toprule
\textbf{Solver} & \textbf{Fit (s)} & \textbf{vs.\ RadVel} & \textbf{$\mu$s per solve} \\
\midrule
Newton  & 0.89 & 17.7$\times$ & 7.0 \\
Danby   & 0.96 & 21.1$\times$ & 7.2 \\
Markley & 1.10 & 23.7$\times$ & 8.4 \\
NWM11   & 1.26 & 27.1$\times$ & 9.3 \\
NWM9    & 1.53 & 32.3$\times$ & 10.9 \\
\bottomrule
\end{tabular}
\end{table}

\paragraph{Real data.} The injection study uses one orbit, so we also
refitted three published datasets with each solver: K2-131 \cite{dai2017}
(70 measurements), HD~164922 \cite{hd164922} (401) and K2-24
\cite{petigura2016,petigura2018} (32). The period, and the transit time
where there is one, were held fixed as in RadVel's own examples, because
left free the sparse data let the fit wander to an alias. The five solvers
agree on the fitted eccentricity to within $6\times10^{-7}$ for K2-131,
$7\times10^{-9}$ for HD~164922 and $7\times10^{-9}$ for K2-24. These
fitted values are not orbit measurements: HD~164922 and K2-24 are
multi-planet systems fitted with one planet. The check only shows that the
solvers agree with each other on real data, as they do on the injected
orbit.

\subsection{Summary of findings}

Table~\ref{tab:synth} collects the results.

\begin{table}[h]
\centering\small
\caption{Findings by question.}
\label{tab:synth}
\begin{tabular}{@{}p{2.3cm}p{5.6cm}@{}}
\toprule
\textbf{Question} & \textbf{Answer} \\
\midrule
Order        & Both papers' claims hold (8.93 and 10.76 measured). \\
Iterations   & With-memory methods need the fewest (1.38 vs.\ 1.94 for Danby). \\
Model cost   & 17\% cheaper than Danby, from the eighth-order parent, not the memory. \\
Clock time   & Slower in Python: 1.4--1.8$\times$ per solve, 1.3--1.7$\times$ per fit. \\
Corner       & Published versions fail 33--37$\times$ as often as Danby. \\
Cause        & The self-accelerating parameter, where $f' \approx 0$. \\
Fix          & Guarded NWM11: 74 $\to$ 3 failures, same order and iterations. \\
Guess        & Napier's guess is best overall; for NWM9 it beats changing method. \\
Downstream   & Solver error is $\ge 1{,}500\times$ below noise. \\
\bottomrule
\end{tabular}
\end{table}

% =====================================================================
\section{Conclusion and Future Work}

We reimplemented two recent with-memory methods, confirmed that they
reach their published orders, and tested them on Kepler's equation against
the methods astronomers already use. Their high order is real: they need
fewer iterations than any other iterative method, and on our cost model
their typical solve is 17\% cheaper than Danby's. That saving comes from
one-iteration solves, where the memory is not yet used, and in pure Python
it does not survive onto the clock.

Their weakness is the corner where $f' \approx 0$. There the
self-accelerating parameter, derived for iterates already near the root,
becomes as large as the derivative it corrects and throws the step away.
Switching the memory off showed that this, not the base method, causes
most failures. A safeguard that keeps the parameter only while its
correction is less than half of $|f'|$ removes almost all of NWM11's
failures and leaves its order and iteration count unchanged. With this
change, NWM11 is as reliable as Danby on our grid and needs fewer
iterations. Separately, a better starting guess is the cheapest
improvement of all: Napier's guess helps every solver on most of the
domain, and for NWM9 it is worth more than changing the method.

For RadVel's working range, Danby or Markley remain the practical choice
in Python. In a compiled kernel, guarded NWM11 with the canonical guess is
the candidate worth testing.

\subsection{Limitations}

\begin{enumerate}
\item \textbf{Python only.} All timings are in CPython, which has no fused
      \texttt{sincos} and makes arithmetic expensive. This works against
      the with-memory methods specifically.
\item \textbf{One machine.} The cost weights were measured on one machine
      and would shift by a few percent on another. Each RadVel fit was
      timed once, so differences below about 0.2\,s are noise.
\item \textbf{The guard constant.} We chose $c = 1/2$ on the same grid we
      report. Values of $0.1$ and $0.01$ gave similar results, but the
      guard has not been tested on a separate set of points.
\item \textbf{NWM9's remaining failures.} The guard does not touch the 17
      failures that come from NWM9's parent method, or its order collapse
      at $e \ge 0.99$.
\item \textbf{One injected orbit.} The downstream study uses one orbit
      with $e = 0.35$. The margin over noise is large, but more eccentric
      orbits were not tested.
\item \textbf{The incumbent.} RadVel's compiled Danby, with its own guess
      and a $10^{-12}$ tolerance, was not run under our cost counting. Our
      ``Danby'' is our Python version at $10^{-14}$.
\end{enumerate}

\subsection{Future work}

The most direct next step is a compiled version. The cost model predicts
that with a fused \texttt{sincos} the 17\% saving becomes real time, and
that prediction can be tested with a C or Cython port of guarded NWM11.
Second, NWM10, the one-parameter method from the same paper as NWM11,
would be the right control for measuring what the second parameter adds;
NWM9 comes from a different paper with a different base method. Third,
the guard constant should be checked on new test points and on other
equations with small derivatives, to see whether the safeguard helps
with-memory methods in general. Finally, injecting more eccentric orbits
would test the downstream result where the solver has the most chance to
matter.

% =====================================================================
\begin{thebibliography}{99}

\bibitem{radvel}
B.~J. Fulton, E.~A. Petigura, S.~Blunt, and E.~Sinukoff.
RadVel: The Radial Velocity Modeling Toolkit.
\emph{Publications of the Astronomical Society of the Pacific},
130(986):044504, 2018. doi:10.1088/1538-3873/aaaaa8.

\bibitem{radvelsrc}
California Planet Search.
RadVel source code, \texttt{src/kepler.c} and \texttt{radvel/kepler.py}.
\url{https://github.com/California-Planet-Search/radvel}.

\bibitem{murray}
C.~D. Murray and S.~F. Dermott.
\emph{Solar System Dynamics}.
Cambridge University Press, 1999.

\bibitem{danby}
J.~M.~A. Danby.
The solution of Kepler's equation, III.
\emph{Celestial Mechanics}, 40(3--4):303--312, 1987.
doi:10.1007/BF01235847.

\bibitem{markley}
F.~L. Markley.
Kepler equation solver.
\emph{Celestial Mechanics and Dynamical Astronomy}, 63(1):101--111, 1995.
doi:10.1007/BF00691917.

\bibitem{nwm9}
S.~K. Mittal, S.~Panday, and L.~J\"antschi.
Enhanced ninth-order memory-based iterative technique for efficiently
solving nonlinear equations.
\emph{Mathematics}, 12(22):3490, 2024. doi:10.3390/math12223490.

\bibitem{nwm11}
S.~K. Mittal, S.~Panday, L.~J\"antschi, and L.~C. Bolundu\c{t}.
Two novel efficient memory-based multi-point iterative methods for solving
nonlinear equations.
\emph{AIMS Mathematics}, 10(3):5421--5443, 2025.
doi:10.3934/math.2025250.

\bibitem{matthies}
G.~Matthies, M.~Salimi, S.~Sharifi, and J.~L. Varona.
An optimal three-point eighth-order iterative method without memory for
solving nonlinear equations with its dynamics.
\emph{Japan Journal of Industrial and Applied Mathematics}, 33:751--766,
2016. arXiv:1602.07026.

\bibitem{solaiman}
O.~S. Solaiman and I.~Hashim.
Optimal eighth-order solver for nonlinear equations with applications in
chemical engineering.
\emph{Intelligent Automation \& Soft Computing}, 27(2):379--390, 2021.

\bibitem{napier}
K.~J. Napier.
Improved initial guesses for numerical solutions of Kepler's equation.
arXiv:2411.15374, 2024.

\bibitem{kipping}
D.~M. Kipping.
Parametrizing the exoplanet eccentricity distribution with the Beta
distribution.
\emph{Monthly Notices of the Royal Astronomical Society}, 434:L51--L55,
2013. doi:10.1093/mnrasl/slt075.

\bibitem{hd164922}
B.~J. Fulton et al.
Three temperate Neptunes orbiting nearby stars.
\emph{The Astrophysical Journal}, 830(1):46, 2016.
doi:10.3847/0004-637X/830/1/46.

\bibitem{dai2017}
F.~Dai et al.
The discovery and mass measurement of a new ultra-short-period planet:
K2-131b.
\emph{The Astronomical Journal}, 154(6):226, 2017.
doi:10.3847/1538-3881/aa9065.

\bibitem{petigura2016}
E.~A. Petigura et al.
Two transiting low density sub-Saturns from K2.
\emph{The Astronomical Journal}, 152(1):28, 2016.

\bibitem{petigura2018}
E.~A. Petigura et al.
Dynamical masses of two sub-Saturns in the K2-24 system.
\emph{The Astronomical Journal}, 155(1):21, 2018.

\bibitem{mpmath}
F.~Johansson et al.
mpmath: a Python library for arbitrary-precision floating-point
arithmetic, version 1.3.0, 2023. \url{https://mpmath.org}.

\end{thebibliography}

\end{document}
